\section{\S 3.2 弱电解质的电离平衡}\label{sect:弱电解质的电离平衡}
    
    \vspace{-\baselineskip}
    \[\cemh{HA <=>[\text{结合}][\text{电离}] H+ + A-}\]

    当\(v_\text{电离}=v_\text{结合}\)时，实现平衡状态——电离平衡

    \textcolor{red}{既然是平衡，所有平衡原理均适用}

    \subsection{电离度}\label{subsec:电离度}

    
        \vspace{-\baselineskip}
        \[
            \alpha=\frac{\text{已电离的电解质分子数}}{\text{该电解质分子总数}}\times \SI{100}{\percent}
        ,\]
        \[
            \alpha=\frac{\text{电离出的离子浓度}}{\text{电解质浓度}}\times \SI{100}{\percent}
        ,\]
        \[
            \alpha=\frac{\cemh{[H+]}}{c}\times \SI{100}{\percent}
        ,\]
        \[
            \cemh{[H+]}=c\cdot\alpha
        .\]

        \subsubsection{意义}\label{subsubsec:电离度的意义}
            \(\alpha\)越大，弱电解质相对越强，越易电离\textcolor{red}{（同温同浓度）}
            
        \subsubsection{影响电离度的因素}\label{subsubsec:影响电离度的因素}

            \vspace{-\baselineskip}
            \[\cemh{HA <=> H+ + A-}\quad\mathcolor{blue}{\increment H>0}\]

            \begin{enumerate}
                \item 温度：\(T\uprightcurvearrow \Rightarrow \alpha\uprightcurvearrow\)
                \ \phantom{\(c\)}\textcolor{red}{升温促进电离（越热越电离）}
                \item 浓度：\(c\downrightcurvedarrow \Rightarrow \alpha\uprightcurvearrow\)
                \ \phantom{\(T\)}\textcolor{red}{稀释促进电离（越稀越电离）}
                
                    \textcolor{red}{\(\therefore\) 弱酸弱碱溶液每稀释10倍，pH改变小于1。}

                \item 外加物质（依情况具体分析）
                \begin{enumerate}
                    \item 同离子效应：加入相同的离子，抑制电离
                    \item 加入能与离子反应的物质，促进电离
                \end{enumerate}
            \end{enumerate}

    \subsection{电离常数}\label{subsec:电离常数}

        电离平衡常数，简称：电离常数，\(K_\text{i}\)表示。
        
        \vspace{0.2\baselineskip}
        \begin{minipage}{0.45\textwidth}
            \centering
            弱酸电离常数，\(K_\text{a}\)
            \vspace{-0.5\baselineskip}
            \[\cemh{HA <=> H+ + A-}\]
            \vspace{-\baselineskip}
            \[K_\text{a}=\frac{\cemh{[H+]}\cdot\cemh{[A-]}}{\cemh{[HA]}}\]
        \end{minipage}
        \hfil
        \begin{minipage}{0.45\textwidth}
            \centering
            弱碱电离常数，\(K_\text{b}\)
            \vspace{-0.5\baselineskip}
            \[\cemh{BOH <=> B+ + OH^-}\]
            \vspace{-\baselineskip}
            \[K_\text{b}=\frac{\cemh{[B+]}\cdot\cemh{[OH^-]}}{\cemh{[BOH]}}\]
        \end{minipage}
        
        \subsubsection{意义}\label{subsubsec:电离常数的意义}

            \(K_\text{i}\)越大，弱电解质电离程度越大，弱酸（或弱碱）相对越强

        \subsubsection{影响因素}\label{subsubsec:影响电离常数的因素}

            只受温度影响，\textcolor{red}{温度升高\(K_\text{i}\)增大}

            \textcolor{blue}{\(K_\text{i}\)受温度影响较小，温度变化不大时，使用常温下的数据。}
            
        \subsubsection{多元弱酸}\label{subsubsec:多元弱酸的电离常数}

            分步电离，分级常数

            \vspace{-\baselineskip}
            \begin{minipage}{0.45\textwidth}
                \[\cemh{H2CO3 <=> H+ + HCO3^-}\]
            \vspace{-\baselineskip}
                \[K_\text{a1}=\frac{\cemh{[H+]}\cdot\cemh{[HCO3^-]}}{\cemh{[H2CO3]}}\]
            \end{minipage}
            \hfil
            \begin{minipage}{0.45\textwidth}
                \[\cemh{HCO3^- <=> H+ + CO3^{2-}}\]
            \vspace{-\baselineskip}
                \[K_\text{a2}=\frac{\cemh{[H+]}\cdot\cemh{[CO3^{2-}]}}{\cemh{[HCO3^-]}}\]
            \end{minipage}

            \(K_\text{a1}>K_\text{a2}>K_\text{a3}\)，一般相差\numrange{1E4}{1E5}

            \textcolor{blue}{\(\therefore\)多元弱酸的酸性主要由第一步电离决定}
        
        \subsubsection{电离常数的应用}\label{subsubsec:电离常数的应用}

        \noindent   1. 比较弱酸、弱碱强弱

            酸性：
            \[\cemh{H2C2O4} > \cemh{H2SO3} > \cemh{HNO2} > \cemh{HF} > \]
            \[\cemh{HCOOH} > \cemh{CH3COOH} > \cemh{H2CO3} > \cemh{HCN}.\]
            
        \noindent   2. 计算弱酸、弱碱pH
        
            \begin{xmp}
                [电离常数计算弱酸、弱碱pH例]
                {}
                []
                []
                计算\(c=\SI{0.1}{\mole\per\litre}\)醋酸溶液的pH和电离度\(\alpha\)
                \begin{center}
                    \begin{tabular}{r@{\quad}c@{\quad}c@{\quad}c@{\quad}c@{\quad}c}
                        & \(\mathbf{\cemh{CH3COOH}} \) & \(\mathbf{\cemh{<=>}} \) & \(\mathbf{\cemh{CH3COO-}} \) & \(\mathbf{+} \) & \(\mathbf{\cemh{H+}} \) \\
                        起始浓度（\si{\mole\per\litre}） & \(\mathbf{0.1}\) && \(\mathbf{0}\) && \(\mathbf{0}\) \\ 
                        平衡浓度（\si{\mole\per\litre}） & \(\mathcolor{blue}{0.1-x}\) && \(\mathcolor{blue}{x}\) && \(\mathcolor{blue}{x}\) \\
                    \end{tabular}
                \end{center}
                \( \displaystyle K_\text{a} = \frac{x^2}{0.1-x}\approx\frac{x^2}{0.1}(x\ll 0.1)\)

                \( \displaystyle \therefore\cemh{[H+]}=\sqrt{0.1K_\text{a}}=\sqrt{0.1\times 1.8\times 10^{-5}}=\SI{1.34E-3}{\mole\per\litre} \)

                \( \therefore\cemh{pH} = 2.87 \) 

                \( \displaystyle \alpha = \frac{\cemh{[H+]}}{c} = \frac{\num{1.34E-3}}{0.1}\times \SI{100}{\percent} = \SI{1.34}{\percent} \)
                \quad \textcolor{red}{“百里挑一”}
            \end{xmp}

            \textcolor{blue}{由以上计算过程可以得出：}
            \( \displaystyle \cemh{[H+]} = \sqrt{c\cdot K_\text{a}}\)，
            \( \displaystyle \alpha = \sqrt{\frac{K_\text{a}}{c}}\)

            \textcolor{red}{\(\therefore\)弱电解质浓度越大，电离度越小}\hfil（前提：\( \displaystyle \frac{c}{K_\text{a}} \ge 400 \)）

        \noindent   3. 辅助分析离子、分子浓度变化
        
            \begin{xmp}
                [电离常数辅助分析离子、分子浓度变化例]
                {}
                []
                []
                往\SI{0.1}{\mole\per\litre}醋酸溶液中\Circled{1}加入少量纯醋酸、\Circled{2}通入少量\cemh{HCl}气体，\( \displaystyle \frac{\cemh{[H+]}}{\cemh{[CH3COOH]}} \)如何变化？

                \( \displaystyle \frac{\cemh{[H+]}}{\cemh{[CH3COOH]}} = \frac{K_\text{a}}{\cemh{[CH3COO^-]}} \)
                \quad\textcolor{red}{\(\therefore\)\Circled{1}变小，\Circled{2}变大}
            \end{xmp}

            请动脑筋算出\SI{25}{\degreeCelsius}时纯水的电离度。

            \vspace{-\baselineskip}
            \[
                \alpha = \frac{\SI{1.0E-7}{\mole\per\litre}\times\SI{1}{\litre}}{\frac{\SI{1000}{\gram}}{\SI{18}{\gram\per\mole}}}\times\SI{100}{\percent} \mathcolor{blue}{= \num{1.8E-9}\ (\SI{1.8E-7}{\percent})}
            .\]

            \textcolor{red}{“亿里挑一”}

        \noindent   4. 比较弱酸根离子结合\cemh{H+}的能力

            \begin{xmp}
                [电离常数比较弱酸根离子结合H+的能力例]
                {}
                []
                []
                \cemh{CH3COO^-}、\cemh{CO3^2-}、\cemh{CN^-}、\cemh{HSO3^-}四种离子中最容易结合\cemh{H+}的是？
                \textcolor{blue}{（所需数据见 \textbf{表~\ref{tab:弱酸-课本3.3}} ）}

                \textcolor{red}{电离常数越小，越难电离，则反向结合（常数越大）越容易}

                \textcolor{blue}{\(\therefore\)结合\cemh{H+}能力：\(\cemh{CO3^2-}>\cemh{CN^-}>\cemh{CH3COO^-}>\cemh{HSO3^-}\)}
            \end{xmp}

        \noindent   5. 判断酸碱反应进行的方向

            \begin{xmp}
                [电离常数判断酸碱反应进行的方向例]
                {}
                []
                []
                根据 \textbf{表~\ref{tab:弱酸-课本3.3}} 的数据（\(K_\text{a}(\cemh{HClO})=\num{2.95E8}\)）分析下列反应是否发生，若发生写出反应的离子方程式。

                \Circled{1}往\cemh{NaCN}溶液中滴入甲酸溶液
                
                \quad 〔\cemh{CN- + HCOOH \xleq{} HCN + HCOO^-}〕

                \Circled{2}往\cemh{Na2CO3}溶液中滴入少量甲酸溶液
                
                \quad 〔\cemh{CO3^2- + HCOOH \xleq{} HCO3^- + HCOO-}〕

                \Circled{3}往\cemh{Na2CO3}溶液中滴入过量甲酸溶液
                
                \quad 〔\cemh{CO3^2- + 2HCOOH \xleq{} H2O + CO2 ^ + 2HCOO-}〕

                \Circled{4}往\cemh{NaCN}溶液中通入过量\cemh{SO2}气体
                
                \quad 〔\cemh{CN- + H2O + SO2 \xleq{} HCN + HSO3^-}〕

                \quad 〔\cemh{CN- + HSO3^- \xleq{} 2HCN + SO3^2-}〕

                \Circled{5}往\cemh{NaCN}溶液中通入过量\cemh{CO2}气体
                
                \quad 〔\cemh{CN- + H2O + CO2 \xleq{} HCN + HCO3^-}〕

                \Circled{6}往\cemh{NaCN}溶液中通入少量\cemh{SO2}气体
                
                \quad 〔\cemh{CN- + H2O + SO2 \xleq{} HCN + HSO3^-}〕

                \Circled{7}往\cemh{NaCN}溶液中通入少量\cemh{CO2}气体
                
                \quad 〔\cemh{CN- + H2O + CO2 \xleq{} HCN + HCO3^-}〕

                \Circled{8}往\cemh{NaClO}溶液中通入过量\cemh{SO2}气体
                
                \quad 〔\cemh{ClO- + H2O + SO2 \xleq{} HClO + HSO3^- \xleq{} Cl- + SO4^2- + 2H+ }〕

                \makebox[0em][r]{〔}\Circled{9}往\cemh{Na2CO3}溶液中通入少量\cemh{SO2}气体\makebox[0em][l]{〕}
                
                \quad 〔\cemh{CO3^2- + H2O + SO2 \xleq{} HCO3^- + HSO3^-}〕
                
                \quad 〔\cemh{CO3^2- + HSO3^- \xleq{} HCO3^- + SO3^2-}〕

                \makebox[0em][r]{〔}\Circled{10}往\cemh{Na2CO3}溶液中通入过量\cemh{SO2}气体\makebox[0em][l]{〕}
                
                \quad 〔\cemh{CO3^2- + H2O + SO2 \xleq{} HCO3^- + HSO3^-}〕
                
                \quad 〔\cemh{HCO3^- + \cancel{\cemh{H2O}} + SO2 \xleq{} \cancel{\cemh{H2O}} + CO2 + SO3^2-}〕

            \end{xmp}

            \textcolor{red}{\Circled{1}强制弱\quad\Circled{2}依据量逐步判断\quad ＊\Circled{3}注意氧还}
